\documentclass[../candidatura.tex]{subfiles}
\graphicspath{{\subfix{../imagenes/}}}
\begin{document}

\section{Conectores}

\avisoesbozo

Para hacer conexiones en sintetizadores modulares, es necesario que nos detengamos en explicar las distintas estrategias, las nomenclaturas, los estándares y los contextos culturales en que ocurren.

Esta lista está ordenada de menor a mayor número de contactos en cada conector, y luego de menor a mayor tamaño en dimensiones físicas.

\subsection{Cables con un contacto}

\subsubsection{Cables banana}

Los cables banana son usados por distintos sintetizadores modulares.

\subsection{Cables con dos contactos}

\subsubsection{Cables TS de 1/8 de pulgada}

Su nombre TS responde a las iniciales en inglés de sus partes Tip (punta) y Sleeve (manga).

En la punta está la señal de interés, y en la manga la señal de tierra.

Estos cables son usados en módulos Eurorack para interconectar módulos, permitiendo que los módulos compartan sus tierras.

\subsubsection{Cables TS de 1/4 de pulgada}

Es la versión de mayor diámetro del cable TS de 1/8 de pulgada.

Son usados en guitarras y bajos eléctricos, y en los dispositivos asociados como efectos y amplificadores.

\subsection{Cables con tres contactos}

\subsubsection{Cables TRS de 1/8 de pulgada}

Su nombre TRS responde a las iniciales en inglés de sus partes Tip (punta), Ring (anillo) y Sleeve (manga).

Estos cables son usados en instrumentos estéreo, para compartir una señal izquierda, otra derecha, y otra de tierra.

\subsubsection{Cables TRS de 1/4 de pulgada}

Es la versión de mayor diámetro del cable TRS de 1/8 de pulgada.

Son usados en distintas configuraciones: en instrumentos estéreo, en configuraciones más complejas del tipo inserto, y también en conexiones mono balanceadas, aunque para estas predomina el tipo XLR.

\subsubsection{Cables XLR}

Los conectores XLR tienen 3 contactos.

Se usan para señales mono balanceadas, en dispositivos como micrófonos, o para hacer conexiones a grandes distancias, sin tener problemas de ruido.

\newpage

\end{document}
